\vssub
\subsection{~\ws\ 建模框架}
\vssub

\ww\ 是一个社区化的波浪建模框架，纳入了风浪建模与动力学领域的最新科学进展。

该框架的核心是 \ws\ 第三代波浪模式。该模式由美国国家环境预报中心
（NOAA/NCEP）开发，其思想承继自 WAM 模式 \citep{bk:WAM94}。
当前框架从早期的 WAVEWATCH I 和 II 模式包
\citep{tol:JPO91b, tol:JPO92} 演进而来，并在控制方程、模式结构、
数值方法和物理参数化等许多重要方面不同于其前身。

\ws\ 求解随机相位谱作用量密度平衡方程，谱以波数--方向形式表示。
该方程的隐含假设是：介质属性（水深和流）以及波场本身的时空变化尺度，
均远大于单个波的变化尺度。模式包含用于浅水（近岸破波带）应用的选项，
也支持网格点的淹没和露出。波谱传播可以在规则网格（直角网格或曲线网格）
和非结构（三角形）网格上求解，也可以将这些网格单独使用或组合为多重网格拼接。

物理过程的源项包括如下参数化：由风作用导致的波浪增长；非线性共振波--波相互作用的精确形式和参数化形式；
由波--海底相互作用引起的散射；三波相互作用；以及由白帽破碎、底摩擦、近岸破碎、
泥沙和海冰相互作用导致的耗散。模式包含多种缓解 Garden Sprinkler Effect 的方法，
并计算其他转换过程，例如表层流对风场和波场的影响，以及未解析岛屿导致的子网格阻挡。

\ws\ 的输入可以通过外部文件提供，也可以通过 OASIS 或 ESMF/NUOPC 框架耦合提供。
输入数据会在波浪模式驱动程序中动态更新，可包括海冰覆盖、泥、流场、
可移动海床上的耗散底部属性，以及供用户进一步开发的数据同化占位模块所需的数据。

\ws\ 采用 ANSI 标准 FORTRAN 90 编写，完全模块化并且全部采用可分配内存。
模式支持传统的单向嵌套，也支持“拼接”或多重网格方法；在后一种方法中，
任意数量的网格都可以同时考虑，并且所有网格之间可以进行完全双向相互作用。
单一网格或多重网格拼接也可以用作移动参考系，从而支持远离海岸的飓风高分辨率建模。

波能谱使用恒定方向增量（覆盖所有方向）和空间变化的波数网格进行离散。
可使用一阶、二阶和三阶精度的数值格式描述波浪传播。源项通过动态调整的时间步进算法进行时间积分，
该算法会在谱快速变化的条件下集中计算量。\ws\ 可以选择编译为包含基于 OpenMP 编译器指令的共享内存并行，
也可以编译为使用消息传递接口（MPI）的分布式内存环境。
